% ============================================================================
% PATCH v5.7.3.04: Units Bridge Clarification
% ============================================================================
% Issue: κ is a free parameter disguised as progress
% Fix: Explicitly label as ansatz, provide measurement protocol, quarantine from theorems
% ============================================================================

\subsection{The Dimensional Bridge: Status Clarification}
\label{sec:units-bridge-status}

\begin{tcolorbox}[colback=red!5!white,colframe=red!75!black,title=Critical Status Update]
The conversion constant $\kappa$ in $\varphi = \kappa \tau_{\mathrm{lat}}$ is currently:
\[
\boxed{\text{[Ansatz / Free Parameter]}}
\]
\textbf{Not} a derived quantity. Any claim depending on $\kappa$ inherits this status.
\end{tcolorbox}

\subsubsection{What We Have vs. What We Need}

\begin{tabular}{p{0.45\textwidth}p{0.45\textwidth}}
\textbf{Currently Established} & \textbf{Required for Derivation Status} \\ \hline
\textbf{Dimensional consistency:} $[\kappa] = L^2 T^{-2} (cy/b)^{-1}$ & \textbf{First-principles derivation:} $\kappa = f(\text{substrate parameters})$ \\
\textbf{Formal definition:} $\kappa = \hbar f_U / (m_P c^2)$ & \textbf{Independent prediction:} Use $\kappa$ to predict new observable \\
\textbf{Consistency check:} Recovers Newtonian limit & \textbf{Measurement protocol:} How to measure $f_U$ independently \\
\end{tabular}

\paragraph{Current Limitations}
\begin{enumerate}
\item \textbf{Circularity risk:} Using $\kappa$ to ``derive'' $G$ then checking against measured $G$ is circular unless $\kappa$ is fixed independently.

\item \textbf{Propagation of uncertainty:} Any downstream claim that depends numerically on $\kappa$ cannot be labeled [Theorem] unless $\kappa$ is derived.

\item \textbf{Calibration leakage:} If $\kappa$ is ultimately set by matching $G$, then claims using $\kappa$ depend on empirical calibration.
\end{enumerate}

\subsubsection{Updated Protocol for κ-Dependent Claims}

\begin{tcolorbox}[colback=yellow!5!white,colframe=yellow!75!black,title=Protocol: Handling κ-Dependent Claims]
For any claim $C$ that depends on $\kappa$:
\begin{enumerate}
\item \textbf{Label:} $C$ must be labeled [Ansatz] or [Conjecture], not [Theorem].
\item \textbf{Dependencies:} Must list ``$\kappa$ ansatz'' as a dependency.
\item \textbf{Falsifiability:} Must provide test independent of $\kappa$ matching.
\item \textbf{If calibrated:} Must be in Track B (Calibrated Models), not Track A (Formal Toy Models).
\end{enumerate}
\end{tcolorbox}

\subsubsection{Measurement Protocol for κ}
To elevate $\kappa$ from ansatz to derived constant, we need:

\begin{mytheorem}{Required Measurement Theorem}
\label{thm:kappa-measurement}
There exists an experimental protocol $\mathcal{P}$ that:
\begin{enumerate}
\item Measures the Universal Clock Frequency $f_U$ independently of gravitational phenomena.
\item Uses this measurement to compute $\kappa = \hbar f_U / (m_P c^2)$.
\item Predicts $G = (8\pi)^{-1} \kappa^2 / c^4$ (or similar relation).
\item This prediction matches the measured value of $G$ within experimental error.
\end{enumerate}
\end{mytheorem}

\paragraph{Concrete Proposal for $\mathcal{P}$}
\begin{enumerate}
\item \textbf{Substrate resonance:} Look for signatures of $f_U$ in:
\begin{itemize}
\item Vacuum fluctuation spectra at $\sim 10^{43}$ Hz
\item Planck-scale discreteness in high-energy collisions
\item Anomalous timing in precision atomic clocks
\end{itemize}

\item \textbf{Independent prediction:} Use $\kappa$ to predict:
\begin{itemize}
\item Variation of $G$ with computational load (testable in dense environments)
\item Quantum gravity corrections to black hole entropy
\item Specific CMB non-Gaussianity pattern from discrete sampling
\end{itemize}

\item \textbf{Falsification condition:} If $\kappa$ measured via $\mathcal{P}$ gives $G$ prediction outside $6.67430 \pm 0.00015 \times 10^{-11} \text{m}^3\text{kg}^{-1}\text{s}^{-2}$, the ansatz is falsified.
\end{enumerate}

\subsubsection{Updated Gravitational Constant Derivation Status}

\begin{tcolorbox}[colback=blue!5!white,colframe=blue!75!black,title=Status: $G$ from $|Co_0|$]
\textbf{Current status:} [Calibration] in Track B.

\textbf{What it gives:} Numeric match $G \approx \frac{c^3}{\hbar} \left(\frac{\hbar}{m_P c}\right)^2 \frac{1}{|Co_0|}$.

\textbf{What it does not give:} First-principles derivation. The match could be numerical coincidence.

\textbf{Required quarantine:} This calibration must \textbf{not} be used to support other claims. Downstream theorems cannot depend on it.
\end{tcolorbox}

\subsubsection{Interim Workaround: Dimensionless Predictions}
Until $\kappa$ is derived, focus on dimensionless predictions that cancel $\kappa$:

\begin{align*}
\frac{\text{DM ratio}}{G} &\quad\text{(still has $G$)} \\
\frac{\text{CMB temperature}}{m_P c^2} &\quad\text{(Planck scale)} \\
\frac{\Lambda}{m_P^2} &\quad\text{(cosmological constant)}
\end{align*}

Or make purely qualitative predictions that don't depend on $\kappa$'s numerical value.

\paragraph{Summary}
The units bridge remains the weakest link in the L-ToEC formalization. Addressing it requires either:
\begin{enumerate}
\item A first-principles derivation of $\kappa$, or
\item An independent measurement protocol for $f_U$, or
\item Reformulating all predictions to be dimensionless.
\end{enumerate}

Until then, all $\kappa$-dependent claims must be clearly labeled as ansatzes.
